\section{Introduction}
\label{sec:intro}

In 1918, after compiling spring phenology records across eastern North America, Andrew Delmar Hopkins proposed a simple rule: spring climbs at four days per four hundred feet of elevation, or 3.28~days per 100~meters \citep{hopkins1918}. A century later, PhenoCam networks in eastern North America \citep{richardson2019} and satellite analyses across mountain gradients in southwest China \citep{jiang2022} confirmed that deciduous canopies broadly follow elevational cooling, while local terrain and microclimate routinely bend this smooth progression.

Elevation alone does not govern how mountain canopies warm or cool. Satellite and field studies have documented nighttime cold-air pooling delaying spring green-up or hastening autumn senescence in mountain valleys \citep{fisher2006, elmore2012, ward2018, pastore2024}. Most recently, \citet{mizuno2026} combined Harmonized Landsat and Sentinel-2 (HLS) phenology with nighttime thermal imagery inside a Japanese mountain basin to show that cold-air pooling can reverse standard elevational phenology. Two open questions remain. First, single-basin studies and local concavity indices such as Topographic Position Index (TPI) fail to separate freely draining gorges from bedrock-rimmed sinks across a replicated sample. Second, dissected tablelands pair microclimate inversions with abrupt lithologic, silvicultural, and forest-type boundaries, so departures from Hopkins' slope can conflate atmospheric pooling, soil hydrology, harvest history, and species turnover.

Topographic closure governs cold-air retention at the base of a karst escarpment. Along the Southern Cumberland Plateau of Tennessee and northern Alabama, open-drainage gorges and enclosed karst blind valleys incise the same Mississippian limestone strata to create a replicated natural contrast. An open gorge lets dense nocturnal air drain freely toward the lowlands, whereas a closed karst sill blocks surface outflow and forces cold air to accumulate over the canopy. At the summit of the escarpment, flat Pennsylvanian sandstone benches host late-leafing upland oak communities adjacent to steeper rocky rims, which bends the upper end of the elevation profile (\cref{fig:hopkins_decomp}).

To separate these mechanisms, we combine offline Priority-Flood depression filling \citep{barnes2014} across a 170~km transect with topographically corrected Sentinel-2 red-edge phenology, BRDF-normalized NASA HLS imagery \citep{claverie2018}, 25-year harmonized Landsat records, public timber-harvest polygons, and cloud-leak-filtered NASA ECOSTRESS thermal scenes. Relative to \citet{mizuno2026}, our contribution is a replicated structural extension: we evaluate 11~forested closed-karst basins matched with adjacent open-drainage gorges, adjust explicitly for relief and elevation imbalances, stratify within identical forest classes, and verify nocturnal cooling across seven basins in tile \texttt{16SED}. We evaluated every hypothesis against internal decision criteria that record partial passes and falsified outcomes alongside supported signals.

Three main findings emerge, each bounded by explicit inferential limits. First, half-amplitude deciduous MODIS green-up averages 3.19~days per 100~meters and matches Hopkins' benchmark, while earlier onset thresholds lapse faster because of phenological stage definition. Second, enclosed karst basins compress the deciduous growing season by 9.1~days after relief adjustment and persist in rim-free floor-vs-floor comparisons inside identical oak and cove guilds. Binary closure outperforms local TPI across our 22~sites, and filtered ECOSTRESS confirms cooling across all seven covered basins plus wind collapse at Lost Cove, though swath-level Manski bounds span zero and sills below 22~meters remain unobserved. Third, flat sandstone plateau benches lag steep rims by 2.6~days on unharvested stands with no response to modeled fragipan soils, a pattern consistent with upland oak guild turnover pending plot-level tree tallies.
